% Part: methods
% Chapter: proofs
% Section: reading-proofs

\documentclass[../../../include/open-logic-section]{subfiles}

\begin{document}

\olfileid{mth}{prf}{rea}
\olsection{阅读证明}

你在教科书和论文中看到的证明，很少会像我们前面的例子那样给出所有细节。
作者通常不会特意说明何时进行了分情况讨论、何时使用了间接证明，
也不会每次都提到用到了哪个定义。
因此，在阅读教科书中的证明时，你往往需要自己补充这些细节才能理解。
这样做也是掌握证明中各种手法的好练习。
让我们来看一个例子。

\begin{prop}[吸收律]
对任意集合 $A$, $B$，
\[
A \cap (A \cup B) = A
\]
\end{prop}

\begin{proof}
如果 $z \in A \cap (A \cup B)$，那么 $z \in A$，所以 $A \cap (A \cup B)
\subseteq A$。现在假设 $z \in A$。那么也有 $z \in A \cup B$，
因此也有 $z \in A \cap (A \cup B)$。
\end{proof}

前面这个吸收律的证明非常浓缩。它没有提到任何用到的定义，没有在正式证明前先声明
“我们要证明……”，等等。我们来拆解它。
被证明的命题是一个关于任意集合 $A$ 和 $B$ 的一般性断言，
当证明中提到 $A$ 或 $B$ 时，它们都是代表任意集合的变元。
证明所建立的一般性断言，正是证明集合相等所需要的内容，
即等式左边的每个元素也都是右边的元素，反之亦然。

\begin{quote}
“如果 $z \in A \cap (A \cup B)$，那么 $z \in A$，所以 $A \cap (A \cup B)
  \subseteq A$。”
\end{quote}

这是证明等式的第一部分：它确立了如果任意~$z$ 是左边的元素，
那么它也是右边的元素，即 $A \cap (A \cup B) \subseteq A$。
假设 $z \in A \cap (A \cup B)$。由于 $z$ 是两个集合交集的元素当且仅当
它是这两个集合各自的元素，我们可以得出结论：$z \in A$ 并且 $z \in A \cup B$。
特别地，$z \in A$，这正是我们想要证明的。
由于这就是第一部分需要完成的全部工作，我们知道证明的其余部分必定是
第二部分的证明，即证明 $A \subseteq A \cap (A \cup B)$。

\begin{quote}
“现在假设 $z \in A$。那么也有 $z \in A \cup B$，并且
因此也有 $z \in A \cap (A \cup B)$。”
\end{quote}

我们从假设 $z \in A$ 开始，因为我们要证明的是：对任意~$z$，
如果 $z \in A$ 那么 $z \in A \cap (A \cup B)$。
要证明 $z \in A \cap (A \cup B)$，我们必须证明（根据“$\cap$”的定义）
 $z \in A$ 并且 $z \in A \cup B$。这里 就是我们的假设，
所以无需再证，因此证明中不再提及它。
对于，回想一下：$z$ 是若干集合之并的元素当且仅当它至少是这些集合中
某一个的元素。由于 $z \in A$，且 $A \cup B$ 是 $A$ 和 $B$ 的并，因此这一条件成立。
所以 $z \in A \cup B$。我们已经证明了 $z \in A$ 和 $z \in A \cup B$，
因此，根据“$\cap$”的定义，$z \in A \cap (A \cup B)$。
证明没有提及这些定义；它假设读者已经掌握了它们。
如果你还没有掌握，你就需要回头复习一下它们的定义。
然后你还需要认识到，为什么从 $z \in A$ 能推出 $z \in A \cup B$，
以及从 $z \in A$ 和 $z \in A \cup B$ 能推出 $z \in A \cap (A \cup B)$。

下面是上面证明的另一个版本，其中所有细节都被显式地给出：
\begin{proof}{}
[根据集合相等的定义，要证明 $A \cap (A \cup B) = A$，我们需要证明
  (a) $A \cap (A \cup B) \subseteq A$ 和 (b) $A \subseteq A \cap (A \cup B)$。
  (a): 根据 $\subseteq$ 的定义，我们需要证明如果 $z \in A \cap (A \cup B)$，
  那么 $z \in A$。] 如果 $z \in A \cap (A \cup B)$，那么 $z \in A$
  [因为根据 $\cap$ 的定义，$z \in A \cap (A \cup B)$ 当且仅当 $z \in A$ 且 $z \in A \cup B$]，
  所以 $A \cap (A \cup B) \subseteq A$。
  [(b): 根据 $\subseteq$ 的定义，我们需要证明如果 $z \in A$，
  那么 $z \in A \cap (A \cup B)$。] 现在假设 [(1)] $z \in A$。
  那么也有 [(2)] $z \in A \cup B$ [因为由 (1) 得 $z \in A$ 或 $z \in B$，
  根据 $\cup$ 的定义，这意味着 $z \in A \cup B$]，
  因此也有 $z \in A \cap (A \cup B)$ [因为 $\cap$ 的定义要求 $z \in A$，即 (1)，
  且 $z \in A \cup B$，即 (2)]。
\end{proof}

\begin{prob}
展开以下 $A \cup (A \cap B) = A$ 的证明，要求你指明所用到的所有推理模式，
说明每一步为何能从之前的假设或已确立的论断推出，以及我们需要在何处诉诸哪个定义。
\begin{proof}
  若 $z \in A \cup (A \cap B)$，则 $z \in A$ 或 $z \in A \cap B$。
  若 $z \in A \cap B$，则 $z \in A$。
  任何 $z \in A$ 也属于 $A \cup (A \cap B)$。
\end{proof}
\end{prob}

\end{document}